\section{The complete transitive $A_7$-cover gate}
\label{sec:a7-cores}

\paragraph{Status.}
New root derivation. The complete proof and all source scopes were
read by the author and an independent reviewer; both checks passed.
The complete GAP spectrum computation, the separate exact Python
reconstruction and an independently written polynomial-Fourier
certificate passed. The independent audit is recorded in
\texttt{A7\_CORE\_PEER\_AUDIT.md}.
The original character tables are identified published inputs, not
newly constructed tables.

\begin{lemma}[D-A7-01: every compatible normalizer]
\label{D-A7-01}
Let $W$ be a nontrivial ordinary irreducible representation of a
quasisimple group with simple quotient $A_7$, and put $a=\dim W$.
If $W$ is not the deleted permutation representation of $A_7$,
every nonscalar finite-order operator normalizing its image satisfies
\[
                  \nu(T):=a-\max_\lambda\dim W_\lambda(T)
                            \ge a/4.                 \tag{A7.1}
\]
Every such representation has $a\ge4$. The only degrees at most six
are four and six, and in either of these nonstandard cases one has
the stronger integral bound $\nu(T)\ge2$.
\end{lemma}

\begin{proof}
The universal perfect cover is $\widetilde A_7=6.A_7$, with cyclic
center of order six. All central quotients are included by using all
ordinary characters of this cover. The complete ordinary-cover list
is stated in \cite[Chapter~4, opening paragraph, p.~31]{BI:HHM};
the ATLAS-derived metadata
\texttt{ExtensionInfoCharacterTable(CharacterTable("A7"))}
are exactly \texttt{["6","2"]}. Also
$\operatorname{Aut}(A_7)=S_7$. The exact official ATLAS table record
\texttt{A7.2} identifies both $S_7$ and
\texttt{AutomorphismGroup(AlternatingGroup(7))}; see
\cite{HD:A7Tables}. These are cited classification inputs.

Use the complete tables \texttt{6.A7} and \texttt{6.A7.2}, with
their original values, class sizes, element orders, power maps and
core fusion. Their input SHA-256 is
\texttt{185c2a90af296395c33ad8c7c5e839fadba5fffaeda3aacac2d99c8d20c906ab}.
Every nonlinear core row and every extension row with irreducible
core restriction is selected. The selection and exact spectra give
\[
\begin{array}{c|r|r|r}
\text{table}&\text{selected rows}&\text{nonscalar spectra}&\min a\\\hline
\texttt{6.A7}&39&1326&4\\
\texttt{6.A7.2}&18&558&6.
\end{array}                                                   \tag{A7.2}
\]
The GAP program checks full row orthogonality and reconstructs every
spectrum. The separate Python checker selects extension rows by the
exact norm of their restricted character, verifies the full fusion,
class-union and power-map compatibility, and recomputes all $1884$
spectra using rational cyclotomic Fourier arithmetic. In more detail,
if $x$ has order $d$, every multiplicity is the integer
\[
             m_j=\frac1d\sum_{r=0}^{d-1}
                    \chi(x^r)\exp(-2\pi i jr/d).          \tag{A7.3}
\]
The Python routine computes its rational field trace with exact
Ramanujan sums; integrality and the sum of all multiplicities are
checked. Every original nonscalar entry is retained, rather than
only entries failing a preliminary character-ratio test.

All these spectra satisfy $4(a-\max m_j)\ge a$, except precisely
extension rows three and four on class~26. Both exceptional rows
have degree six and restrict to core row two; their spectra are
$(5,1)$ and $(1,5)$ up to the choice of initial eigenvalue. Core
row two is the unique degree-six row trivial on all six central
classes. It is the deleted permutation representation: the natural
$A_7$ action is doubly transitive, so its deleted module is
irreducible of degree six. The complete core table gives no other
center-trivial row of that degree. The two degree-four rows have
minimum projective support two, and every nonstandard degree-six
row has minimum support at least two. None of them is the
irreducible restriction of an extension row. The only nonlinear
core degrees at most six are four and six. These assertions are
explicitly tested in \texttt{research/a7\_cores/check.py} and
recorded for every row in \texttt{checked.json}.

It remains to justify that these two tables cover every compatible
linear normalizer, including scalar twists and other abstract
extensions. An automorphism of a represented central quotient
lifts uniquely to its universal perfect cover. An automorphism
of that cover which is trivial on $A_7$ differs from the identity
by a homomorphism from a perfect group to its center; it is therefore
the identity. If the induced automorphism is inner, Schur's lemma
makes $T$ a scalar multiple of a represented core element.

Otherwise adjust by an inner operator to the automorphism induced
by a fixed lift $t$ in the chosen group $6.A_7.2$. A compatible
intertwiner $A$ has $A^2$ equal to $\rho(t^2)$ times a scalar.
Rescale $A$ to remove that scalar. This gives an actual extension
of $W$ to the chosen group, with irreducible core restriction.
The original $T$ is a scalar times an element of that extension.
This argument does not require the chosen extension to split and
also handles an isoclinic extension with a different scalar square.
Projective support is unchanged by the scalar. Thus the complete
table inequalities give exactly (A7.1) and its stated endpoints.
\end{proof}

\begin{lemma}[D-A7-02: the standard tensor endpoint]
\label{D-A7-02}
In the actual sixfold-source setting, an inducing block with
irreducibly acting normal core $A_7^t$ on
$\operatorname{Std}_7^{\otimes t}$, $t\ge2$, cannot have a
nine-dimensional period image. This permits all compatible outer
operators, scalars and permutations of the factors, and a larger
inducing character field.
\end{lemma}

\begin{proof}
Put $e=6^t$. The full normalizer, by $\operatorname{Aut}(A_7)=S_7$
and Schur's lemma on each factor, consists of scalar multiples
$g=\lambda T$ of the rational orthogonal tensor-wreath operators
$T\in\mathcal R_{7,t}$ defined in D-AST-01. That lemma was proved
for all $n\ge3$, and gives, for every integer $j$,
\[
                 \operatorname{tr}(T^j)\ge-e/6.
                                                               \tag{A7.4}
\]
Its elementary proof contracts the tensor indices around each
factor cycle; every resulting standard trace lies in
$\{-1,0,1,\ldots,6\}$. A negative product therefore has absolute
value at most $6^{t-1}$. This also applies to every power.

For any nonreal eigenvalue $\zeta$ of $T$ with multiplicity $m$,
its conjugate has the same multiplicity. Some power has
$\operatorname{Re}(\zeta^j)\le-1/2$. The trace of that power is
at most $e-3m$, so (A7.4) gives $m\le7e/18$.
The tensor standard form has multiplier $\mu(g)=\lambda^2$.
Consequently an element of nontrivial multiplier has moving rank
at least $11e/18$, because its fixed space is a nonreal eigenspace
of this particular rational $T$.

Every other genuine inertia has moving rank at least $e/6$.
A nonscalar factor matrix has projective support at least one
in degree six; convolution of normalized spectra preserves this
lower fraction on a factor-preserving tensor. A nontrivial factor
cycle has the stronger quarter bound, by the weighted basis-word
argument in D-SPOR-03. A nonidentity scalar moves all of $V$.

The multiplier character is nontrivial: otherwise the finite
inducing image would preserve a nondegenerate bilinear form,
have real character, and could not have a character field
containing the original imaginary quadratic field. This is the
field inclusion from D-BI-01, without assuming field equality.
The generating product-one branch tuple has at least two entries
with nontrivial multiplier.

The actual intermediate-cover Euler identity and period
factorization give a genus-zero base, at least twelve genuine
inertias and total moving rank $2e+6$. Genus at least two is
impossible for $e\ge36$; genus one has actual variation rank at
most six. Hence
\[
        2e+6\ge2\frac{11e}{18}+10\frac e6
                  =\frac{26e}{9},
                \qquad e\le\frac{27}{4}<36,             \tag{A7.5}
\]
contradicting $e=6^t\ge36$. This extends the already written
standard-tensor argument to its missing $n=7$ endpoint; it does
not assert that a nonreal-eigenspace bound survives scalar
convolution without first extracting the actual scalar phase.
\end{proof}

\begin{theorem}[D-A7-03: no transitive $A_7$-cover core in the nonsplit gate]
\label{D-A7-03}
An inducing block in the nonsplit Weil-sixfold problem cannot
have a nine-dimensional actual period image if its normal
quasisimple central-product core acts irreducibly, all simple
quotients are $A_7$, and the inducing group is transitive on those
factors. All central covers, compatible outer and scalar actions,
factor permutations and larger inducing character fields are
allowed.
\end{theorem}

\begin{proof}
Write the representation of the lifted core as
$V=\bigotimes_{i=1}^tW_i$, with $e=\prod_i\dim W_i$.
Transitivity implies that all factor characters agree up to
automorphisms of the universal cover. In particular, if one is
the ordinary deleted representation, then all are: this character
is invariant under $\operatorname{Aut}(A_7)$, and its central
kernel is intrinsic. For $t\ge2$, D-A7-02 excludes this case.
For $t=1$, the full normalizer is
$\mathbf C^*S_7$ in its deleted representation. D-BI-08 proves
that rank nine then forces the actual polarization to be split,
contrary to the present hypothesis. Its proof explicitly allows
the inducing character field to be larger initially.

Otherwise every factor is nonstandard. By D-A7-01, every
nonscalar factor-preserving normalizer on $V$ has projective
support at least $e/4$; normalized eigenvalue multiplicities
convolve. A nontrivial factor permutation has the same quarter
bound, and a nonidentity scalar moves all $e$ dimensions.
Since $e\ge4$, the same actual intermediate-cover formula forces
genus zero for rank nine, at least twelve branches and total
moving rank $2e+6$. Thus
\[
                  12\lceil e/4\rceil\le2e+6,
                              \qquad e\le6.             \tag{A7.6}
\]
If $t\ge2$, then $e\ge4^2>6$. If $t=1$, the only degrees
left are four and six. Each genuine nonscalar inertia in either
case has moving rank at least two by D-A7-01; a genuine scalar
has still larger moving rank. Thus $24\le2e+6\le18$, another
contradiction. These exhaust the transitive core cases.
\end{proof}

\paragraph{Geometry, framework and boundary.}
This removes the residual $A_7$ and its double, triple and sixfold
covers from the transitive normal-core gate forced by D-BI-19.
It does not assert a result for an arbitrary nontransitive product
of differently labelled factors. The argument is about actual
finite covers and their weight-one multiplicity variations; the
integer spectrum and tensor identities belong to the associated
representation ring. The exact scheme-spreading, SGA finite-cover,
lambda-ring, K-theory and etale-weight placements are those in the
framework section. It narrows the supplier-existence part of
condition A, cote16 scan050, without producing a cycle or changing
condition I, cote16 scan049 and cote44 Theorem4, or the
Cayley--Hamilton projector in cote16 scans048--049. No new
motivic $L$-value estimate or general negative Hodge-conjecture
conclusion follows. All existing detector constants are unchanged.
